\section{近似闭合：最优核、纤维直径与复杂度--误差前沿}
\label{app:approximate-theory}

精确交换式在估计误差下不稳定；本附录把交换式缺陷作为近似对象，并给出主判据的闭合项。纤维半径--直径关系把含未知宏观核的最优闭合误差化为可由分区直接计算的预测直径。所有存在性结论都来自有限维紧性，不应被解释为高效发现算法。

\subsection{固定连接通道的最优宏观核}

设 $K\in\Stoch(X,X)$、$L\in\Stoch(Y,Y)$，且
$Q\in\Stoch(X,Y)$。定义
\begin{equation}
  \partial_{K,L}(Q):=d_\infty(KQ,QL).
  \label{eq:app-factor-defect}
\end{equation}
当该量不超过 $\eps$ 时，称 $Q$ 为从 $(X,K)$ 到 $(Y,L)$ 的
$\eps$-近似动力学因子。这里 $K,L$ 仍是严格归一化的 Markov 核；近似的
只是交换关系。

固定 $K,Q$，令
\begin{equation}
  \eta_Q(K)
  :=\min_{L\in\Stoch(Y,Y)}\partial_{K,L}(Q),
  \qquad
  \operatorname{Opt}_Q(K)
  :=\argmin_{L\in\Stoch(Y,Y)}\partial_{K,L}(Q).
  \label{eq:app-eta-fixed-q}
\end{equation}

\begin{theorem}[有限空间上的最优核达到]
\label{thm:app-optimal-kernel}
$\operatorname{Opt}_Q(K)$ 非空、紧且凸。因此
式~\eqref{eq:app-eta-fixed-q} 中写的是最小值而非下确界。最优核一般不唯一。
\end{theorem}

\begin{proof}
$\Stoch(Y,Y)$ 是有限维欧氏空间中的非空紧凸多面体。映射
$L\mapsto QL$ 仿射，而
\[
 L\longmapsto d_\infty(KQ,QL)
 =\max_{x\in X}\frac12\sum_{y\in Y}
   |(KQ)(x,y)-(QL)(x,y)|
\]
是连续凸函数，故在紧集上达到最小值。最优集是紧集的闭子集。若
$L_1,L_2$ 均最优，则目标的凸性说明它们的任意凸组合目标值不超过最小值，
因而仍等于最小值。
\end{proof}

随机 $Q$ 下，即使 $\eta_Q(K)=0$，宏观核也未必唯一，因为线性映射
$L\mapsto QL$ 可能不单射。下节的确定性满射情形在零误差时没有这一歧义。

\subsection{确定性分区的 Chebyshev 中心公式}

令 $q:X\twoheadrightarrow Y$，并记
\[
  p_x^q:=(KQ_q)(x,\cdot)\in\Delta(Y),
  \qquad
  p_x^q(y)=K\bigl(x,q^{-1}(y)\bigr).
\]
简记 $\eta_q(K):=\eta_{Q_q}(K)$，并定义纤维预测直径
\begin{equation}
  D_q(K):=\max_{q(x)=q(x')}
  \|p_x^q-p_{x'}^q\|_{\TV}.
  \label{eq:app-fiber-diameter}
\end{equation}

\begin{theorem}[纤维 Chebyshev 中心]
\label{thm:app-chebyshev-center}
有
\begin{equation}
 \eta_q(K)
 =\max_{y\in Y}
   \min_{p\in\Delta(Y)}
   \max_{x\in q^{-1}(y)}\|p_x^q-p\|_{\TV},
 \label{eq:app-chebyshev-formula}
\end{equation}
并且
\begin{equation}
  \boxed{\frac12D_q(K)\le\eta_q(K)\le D_q(K).}
  \label{eq:app-diameter-radius}
\end{equation}
此外，
\begin{equation}
 \eta_q(K)=0
 \quad\Longleftrightarrow\quad
 D_q(K)=0
 \quad\Longleftrightarrow\quad
 q\text{ 对 }K\text{ strong-lumpable}.
 \label{eq:app-zero-approx-exact}
\end{equation}
此时诱导宏观核唯一。
\end{theorem}

\begin{proof}
因为 $(Q_qL)(x,\cdot)=L(q(x),\cdot)$，
\[
 \partial_{K,L}(Q_q)
 =\max_{y\in Y}\max_{x\in q^{-1}(y)}
   \|p_x^q-L(y,\cdot)\|_{\TV}.
\]
$L$ 的不同行可以彼此独立地在 $\Delta(Y)$ 中选择；每个单纯形紧致，故每个
纤维的最小包围半径达到。先逐纤维选中心再取最坏纤维，得到
式~\eqref{eq:app-chebyshev-formula}。

对同一纤维中的 $x,x'$ 与任意中心 $p$，三角不等式给出
\[
 \|p_x^q-p_{x'}^q\|_{\TV}
 \le\|p_x^q-p\|_{\TV}+\|p-p_{x'}^q\|_{\TV}.
\]
所以每个包围半径至少是相应直径的一半，得到左侧不等式。反过来，在每个
纤维任选代表 $x_y$ 并取中心 $p_{x_y}^q$，该纤维内全部距离不超过全局直径，
得到右侧不等式。半径为零当且仅当每个纤维内所有 $p_x^q$ 相同；由
定理~\ref{thm:app-four-equivalences}，这正是 strong lumpability。共同的
纤维分布逐行唯一确定宏观核。
\end{proof}

界~\eqref{eq:app-diameter-radius} 的常数不能仅凭度量空间三角不等式改进；
本文只使用它建立误差为零和误差趋零的等价，不声称 $D_q/2$ 总是精确半径。

\subsection{动力学族中的量词顺序}

若同一表示 $q$ 要服务于一族动力学 $\cK\subseteq\Stoch(X,X)$，本文采用
\begin{equation}
 \boxed{
 \eta_q(\cK)
 :=\sup_{K\in\cK}\eta_q(K)
 =\sup_{K\in\cK}\min_{L_K\in\Stoch(Y,Y)}
 d_\infty(KQ_q,Q_qL_K).}
 \label{eq:app-sup-min}
\end{equation}
这里同一个 $q$ 必须适用于整族动力学，但宏观核 $L_K$ 可以随 $K$ 改变。
这是\emph{$\sup_K\min_L$}，不是共享单一核所要求的
\begin{equation}
 \eta_q^{\mathrm{shared}}(\cK)
 :=\min_{L\in\Stoch(Y,Y)}\sup_{K\in\cK}
 d_\infty(KQ_q,Q_qL).
 \label{eq:app-min-sup}
\end{equation}
紧性与目标函数的下半连续性保证式~\eqref{eq:app-min-sup} 的最小值也达到，
且一般只有
\begin{equation}
  \sup_K\min_L d_\infty(KQ_q,Q_qL)
  \le
  \min_L\sup_K d_\infty(KQ_q,Q_qL).
  \label{eq:app-minimax-order}
\end{equation}
除非另行证明二者相等，正文的系统族判据不产生一个对所有 $K$ 通用的转移核。

定义
\begin{equation}
  D_q(\cK):=\sup_{K\in\cK}D_q(K).
  \label{eq:app-family-diameter}
\end{equation}
若 $q=q_\pi$ 来自分区 $\pi$，正文的记号就是
$\eta_\pi(\cK):=\eta_{q_\pi}(\cK)$ 与
$D_{\cK}(\pi):=D_{q_\pi}(\cK)$。
逐 $K$ 应用定理~\ref{thm:app-chebyshev-center} 后取上确界，得到
\begin{equation}
  \frac12D_q(\cK)\le\eta_q(\cK)\le D_q(\cK).
  \label{eq:app-family-diameter-radius}
\end{equation}
因此 $D_q(\cK)\to0$ 与本文所用
$\sup_K\min_L$ 缺陷趋零等价到固定常数，而不需要先估计任何 $L_K$。

\subsection{核族的任务相对复杂度--误差前沿}

设 $g:X\twoheadrightarrow O$ 非恒定，
$\cK\subseteq\Stoch(X,X)$ 非空。称分区 $\pi$ 任务保真，如果同一
$\pi$-块中的状态具有相同 $g$ 值；等价地，其商映射
$q_\pi:X\twoheadrightarrow X/\pi$ 允许唯一读出 $r_\pi$ 使
$g=r_\pi\circ q_\pi$。对 $|O|\le m\le|X|$，令
\[
 \mathfrak P_m(g):=
 \left\{\pi:\pi\text{ 任务保真且 }|\pi|\le m\right\}.
\]
定义复杂度--误差前沿与直径前沿
\begin{align}
  \mathcal E_{\cK,g}(m)
  &:=\min_{\pi\in\mathfrak P_m(g)}\eta_{q_\pi}(\cK),
  \label{eq:app-complexity-error-curve}\\
  \mathcal D_{\cK,g}(m)
  &:=\min_{\pi\in\mathfrak P_m(g)}D_{q_\pi}(\cK).
  \label{eq:app-diameter-frontier}
\end{align}

\begin{theorem}[核族复杂度--误差前沿的有限达到]
\label{thm:app-complexity-error-curve}
对每个允许的 $m$，式~\eqref{eq:app-complexity-error-curve} 与
式~\eqref{eq:app-diameter-frontier} 的外层最小值都由任务保真分区达到；
对每个固定分区与每个固定 $K\in\cK$，内层最优宏观 kernel 也达到。
并且：
\begin{enumerate}
  \item $m\mapsto\mathcal E_{\cK,g}(m)$ 单调不增；
  \item $\mathcal E_{\cK,g}(|X|)=0$；
  \item $\mathcal E_{\cK,g}(m)=0$ 当且仅当存在一个至多 $m$ 状态、
  对全部 $K\in\cK$ 精确闭合的任务保真商；
  \item
  \begin{equation}
  c_{\cK}^{0}(g)
  =\min\{m:\mathcal E_{\cK,g}(m)=0\}.
  \label{eq:app-exact-frontier-threshold}
  \end{equation}
\end{enumerate}
这里没有声称 $\sup_{K\in\cK}$ 必由某个 $K$ 达到。
\end{theorem}

\begin{proof}
有限集只有有限多个分区，$\mathfrak P_m(g)$ 非空且有限；例如离散分区在
$m=|X|$ 时属于其中。对每个固定分区与固定 $K$，
定理~\ref{thm:app-optimal-kernel} 给出最优宏观核。虽然核族可以无限，
$\eta_{q_\pi}(\cK)$ 对每个分区仍是一个良定义的上确界；外层只在有限个
分区上取最小值，故达到。增大 $m$ 只扩张可行集，得到单调性。离散分区
以 $L_K=K$ 对每个 $K$ 实现零误差。

$\eta_{q_\pi}(\cK)=0$ 当且仅当
$\eta_{q_\pi}(K)=0$ 对全部 $K\in\cK$ 成立；由
式~\eqref{eq:app-zero-approx-exact}，这等价于 $q_\pi$ 对整个核族共同
精确闭合。再结合定理~\ref{thm:app-minimal-task-quotient} 与
式~\eqref{eq:app-exact-complexity-min}，得到后两项。
\end{proof}

对误差容限 $\tau\ge0$，定义
\begin{equation}
 c_{\cK}^\tau(g)
 :=\min\{ |\pi|:\pi\text{ 任务保真且 }
                    \eta_{q_\pi}(\cK)\le\tau\}.
 \label{eq:app-tolerance-complexity}
\end{equation}
则
\begin{equation}
 c_{\cK}^\tau(g)
 =\min\{m:\mathcal E_{\cK,g}(m)\le\tau\},
 \label{eq:app-tolerance-frontier-inverse}
\end{equation}
并且
\begin{equation}
 |O|\le c_{\cK}^\tau(g)\le|X|.
\end{equation}
$c_{\cK}^\tau(g)$ 随 $\tau$ 单调不增；在 $\tau=0$ 时，它与
附录~\ref{app:exact-theory} 构造的最粗精确任务商状态数一致。

逐分区使用式~\eqref{eq:app-family-diameter-radius} 并取最小值，得到
\begin{equation}
 \frac12\mathcal D_{\cK,g}(m)
 \le\mathcal E_{\cK,g}(m)
 \le\mathcal D_{\cK,g}(m).
 \label{eq:app-frontier-radius-diameter}
\end{equation}
因此两条前沿的零点与渐近趋零行为相同。精确情形由
定理~\ref{thm:app-minimal-task-quotient} 给出唯一最粗商；当 $\tau>0$ 时，
达到同一最优状态数或同一前沿值的分区一般不可比较，故不应声称存在唯一
最粗的近似最优分区。规范对象是前沿及各点的全部最优解集合。

当 $\cK=\{K\}$ 时，简记
\[
 \mathcal E_{K,g}:=\mathcal E_{\{K\},g},
 \qquad
 c_K^\tau(g):=c_{\{K\}}^\tau(g).
\]

\subsection{估计核下的稳健性}

\begin{proposition}[固定通道的模型误差转移]
\label{prop:app-model-perturbation}
对同型核 $K,\widehat K\in\Stoch(X,X)$、
$L,\widehat L\in\Stoch(Y,Y)$ 和固定 $Q\in\Stoch(X,Y)$，
\begin{align}
 \left|\partial_{K,L}(Q)-
 \partial_{\widehat K,\widehat L}(Q)\right|
 &\le \alpha(Q)d_\infty(K,\widehat K)
       +d_\infty(L,\widehat L),
 \label{eq:app-defect-perturbation}\\
 |\eta_Q(K)-\eta_Q(\widehat K)|
 &\le \alpha(Q)d_\infty(K,\widehat K).
 \label{eq:app-eta-perturbation}
\end{align}
\end{proposition}

\begin{proof}
度量的反三角不等式、引理~\ref{lem:app-tv-contraction} 分别给出
\begin{align*}
 &\left|d_\infty(KQ,QL)-d_\infty(\widehat KQ,Q\widehat L)\right|\\
 &\qquad\le d_\infty(KQ,\widehat KQ)+d_\infty(QL,Q\widehat L)\\
 &\qquad\le \alpha(Q)d_\infty(K,\widehat K)+d_\infty(L,\widehat L).
\end{align*}
若第二个核固定为同一 $L$，两个关于 $L$ 的优化目标在整个紧可行集上一致
相差至多第一项；分别代入对方的最优解即得
式~\eqref{eq:app-eta-perturbation}。
\end{proof}

\begin{corollary}[确定性分区与复杂度曲线的稳健性]
\label{cor:app-diameter-curve-robustness}
若 $d_\infty(K,\widehat K)\le\delta$，则对每个满射分区 $q$，
\begin{align}
 |\eta_q(K)-\eta_q(\widehat K)|&\le\delta,
 \label{eq:app-eta-q-robust}\\
 |D_q(K)-D_q(\widehat K)|&\le2\delta,
 \label{eq:app-diameter-robust}\\
 |\mathcal E_{K,g}(m)-\mathcal E_{\widehat K,g}(m)|&\le\delta.
 \label{eq:app-curve-robust}
\end{align}
因而
\begin{equation}
 c_K^{\tau+\delta}(g)\le c_{\widehat K}^{\tau}(g),
 \qquad
 c_{\widehat K}^{\tau+\delta}(g)\le c_K^\tau(g).
 \label{eq:app-horizontal-robustness}
\end{equation}
\end{corollary}

\begin{proof}
非恒定任务保真 $q$ 满足 $\alpha(Q_q)=1$，而常值情形只会给出更强界；
式~\eqref{eq:app-eta-q-robust} 来自
命题~\ref{prop:app-model-perturbation}。每个投影行在 $K$ 与
$\widehat K$ 下相距至多 $\delta$；对任意一对行使用反三角不等式，得其
两两距离变化至多 $2\delta$，取最大值得到
式~\eqref{eq:app-diameter-robust}。对每个同一分区的 $\eta$ 值使用第一式，
再分别代入对方曲线的最优分区，得到式~\eqref{eq:app-curve-robust}。若某分区
在 $\widehat K$ 下误差不超过 $\tau$，则它在 $K$ 下误差不超过
$\tau+\delta$；交换二者即得式~\eqref{eq:app-horizontal-robustness}。
\end{proof}

最后给出正文 $\Gamma$ 判据所需的直接稳健形式。设同一有限操作性微观空间
$B$ 上有同索引核族 $\{K_\theta\}_{\theta\in\Theta}$ 与估计族
$\{\widehat K_\theta\}_{\theta\in\Theta}$，且
\[
 \sup_{\theta\in\Theta}d_\infty(K_\theta,\widehat K_\theta)
 \le\delta.
\]
对任务保真分区 $\pi$，沿用
$D_{\cK}(\pi)=\sup_\theta D_{q_\pi}(K_\theta)$，并定义
\begin{equation}
 \Gamma(B,\cK,g)
 :=\min_{\pi\text{ 任务保真}}
 \max\left\{\frac{|\pi|}{|B|},D_{\cK}(\pi)\right\}.
 \label{eq:app-gamma-definition}
\end{equation}
式~\eqref{eq:app-diameter-robust} 对 $\theta$ 一致成立，故
\begin{equation}
 \left|\Gamma(B,\cK,g)-
 \Gamma(B,\widehat{\cK},g)\right|\le2\delta.
 \label{eq:app-gamma-robustness}
\end{equation}
这是确定性的模型扰动结论。只有在另行给出
$d_\infty(K_\theta,\widehat K_\theta)$ 的统计置信界后，它才成为高概率
统计证书；本文不从有限性本身推出样本复杂度。
